\documentclass[11pt]{article}
\usepackage[a4paper,margin=22mm]{geometry}
\usepackage[T1]{fontenc}
\usepackage{lmodern}
\usepackage{amsmath}
\usepackage{microtype}
\usepackage{graphicx}
\usepackage{xcolor}
\usepackage{float}
\usepackage{fancyhdr}
\usepackage{booktabs}
\usepackage[hidelinks]{hyperref}

\definecolor{Accent}{HTML}{21918C}
\definecolor{Muted}{HTML}{5B6570}
\definecolor{Warn}{HTML}{B03A2E}

\pagestyle{fancy}\fancyhf{}
\lhead{\small\color{Muted}Local-score null calibration (confirmatory)}
\rhead{\small\color{Muted}module\_localscore\_crosscheck}
\cfoot{\small\thepage}
\renewcommand{\headrulewidth}{0.4pt}
\setlength{\parskip}{0.5em}\setlength{\parindent}{0pt}

\title{\vspace{-3em}\textbf{\color{Accent}Is BayPass's local-score method calibrated under the null?}\\[0.2em]
\large A pre-specified, structured-vs-unstructured full-SNP test, and whether an
LD-aggregated large-replicate design is needed to show it}
\author{}
\date{\small Generated \today}

\begin{document}
\maketitle
\vspace{-2em}

\section*{\color{Accent}Motivation and question}
BayPass ships its own method for turning per-SNP association evidence into outlier
\emph{regions}: \texttt{compute.local.scores()} (\texttt{baypass\_public/utils/baypass\_utils.R}),
implementing the local-score approach of Fariello et al.\ (2017) and Bonhomme et al.\ (2019). Its
appeal is that it reports region-level, not SNP-level, significance directly from a single
observed scan, without requiring a separate null-calibration step. This document asks two
questions, pre-specified before running any BayPass job:

\begin{enumerate}
\item \textbf{Is the method's own significance threshold trustworthy?} If a covariate with
  \emph{no} real biological signal -- either fully independent of population structure
  (``unstructured'') or built to respect it the same way the pipeline's real covariates do
  (``structured'', drawn as $\mathrm{MVN}(0,\Omega)$) -- still produces many ``significant''
  local-score windows, then the raw \texttt{significant.windows} output cannot be trusted at face
  value, regardless of which real covariate is being tested.
\item \textbf{Is a more elaborate design necessary to show this?} A separate, more expensive
  strand of this pipeline generates very large numbers of null replicates ($10^3$--$10^4$) by
  working at LD-aggregated (Stage-1/Stage-2 cluster) marker resolution, cheap enough to run at
  that scale. This document asks whether that machinery is actually \emph{necessary} to
  demonstrate poor null calibration, or whether a small, full-resolution, directly interpretable
  design already settles the question.
\end{enumerate}

\section*{\color{Accent}Design}
Ten replicate null covariates were drawn as $\mathrm{MVN}(0,\Omega)$ (``structured'', the same
mechanism used throughout this pipeline's real covariates), using the authoritative population
covariance matrix already estimated for this dataset
(\url{module_manuscript_rho05/baypass_stage1/aland_excluded/omega_mat_omega.out}, 19
populations, Åland excluded). Each structured draw was paired with an \emph{unstructured}
counterpart: the same 19 realized values, randomly permuted across populations. Pairing this way
holds the value distribution exactly fixed (identical mean, SD, and value multiset) and isolates
one variable -- whether the covariate's pattern correlates with $\Omega$ -- as the only difference
between the two null types (full derivation in \texttt{config/null\_covariate\_design.md}).

Each of the 10 structured/unstructured pairs was run through BayPass on the full genome-wide SNP
set (1{,}114{,}423 SNPs, \texttt{u\_DIEM.geno}), in two modes mirroring the pipeline's two real
covariate types: a continuous covariate (IS covariate mode, \texttt{-efile}, BF in deciban units)
and a binary mitotype-style contrast (\texttt{-contrastfile}, $C2$/$-\log_{10}(p)$) with the same
7-vs-12 population split as the real mitotype contrast. This gives $10\times2\times2=40$ full-SNP
BayPass runs, none reusing any output, seed, or draw from the earlier exploratory investigation of
this question.

\texttt{compute.local.scores()} was then run independently on all 40 outputs with the same
settings used throughout this pipeline ($\xi=1$, local-score significance threshold $0.01$,
$\mathrm{MAF}>0.2$, minimum 10{,}000 SNPs per chromosome to be included) -- 1{,}083{,}000--1{,}083{,}700
of the 1{,}114{,}423 SNPs were retained per run after the internal MAF filter, consistent across all
40 runs.

\section*{\color{Accent}Result 1: null covariates already produce many \emph{significant} windows}

\begin{table}[H]
\centering\small
\begin{tabular}{lrrrr}
\toprule
 & mean & sd & min & max \\
\midrule
continuous, structured   & 4.6   & 3.47  & 0  & 10 \\
continuous, unstructured & 4.5   & 2.72  & 1  & 8 \\
mitoC2, structured        & 112.9 & 18.23 & 86 & 147 \\
mitoC2, unstructured      & 99.9  & 13.75 & 74 & 122 \\
\bottomrule
\end{tabular}
\caption{Significant local-score windows per replicate, genome-wide, across the 10 replicates of
each null type/mode.}
\end{table}

\begin{figure}[H]
\centering
\includegraphics[width=0.9\textwidth]{../Figures/window_counts_structured_vs_unstructured.pdf}
\caption{Significant window counts per replicate. Lines connect each replicate's paired
structured/unstructured draw. \emph{Note: axis scales differ between panels by design (free
y-scales) -- do not compare the two panels' vertical positions directly; only the within-panel
structured-vs-unstructured contrast and the numbers in the table above are informative. Colour
distinguishes null type; it is not itself evidence -- see numerical summaries.}}
\end{figure}

In continuous mode, a covariate with no real signal produces a modest number of false windows
(under 5, on average). In mitoC2 (contrast) mode, a covariate with no real signal produces
\textbf{over 100 ``significant'' windows per replicate on average} -- roughly 20--25$\times$ more
than the continuous mode's false-positive rate, from an identical significance threshold
($\mathrm{local\text{-}score\ } p<0.01$) applied to an equally fabricated null. This alone
demonstrates that \texttt{compute.local.scores()}'s raw output is not usable as a trustworthy
significance call without an explicit null-based correction: a method that would be expected to
report on the order of a handful of false regions genome-wide under a well-calibrated $p<0.01$
threshold is instead reporting well over a hundred, for this contrast statistic on this dataset.

The structured-vs-unstructured contrast within each mode is smaller and, at this replicate count,
not conclusive:

\begin{table}[H]
\centering\small
\begin{tabular}{lrrr}
\toprule
mode & mean difference (structured $-$ unstructured) & paired $t$-test $p$ & paired Wilcoxon $p$ \\
\midrule
continuous & $+0.1$ windows  & 0.946 & 1.000 \\
mitoC2     & $+13.0$ windows & 0.078 & 0.064 \\
\bottomrule
\end{tabular}
\end{table}

Continuous mode shows no detectable difference between the two null types. mitoC2 mode shows a
directional excess for the structured null (higher in 8 of 10 paired replicates) that is
suggestive but does not reach conventional significance at $n=10$ ($p\approx 0.06$--$0.08$) --
reported here as an honest non-significant result, not as evidence either way for whether
$\Omega$-linked structure specifically (as opposed to an unstructured covariate of the same
magnitude) inflates local-score detection. What \emph{is} unambiguous is that \emph{both} null
types, structured and unstructured alike, already produce a large number of false windows in
mitoC2 mode -- the calibration failure is not specific to population structure.

\section*{\color{Accent}Result 2: a third to half of observed ``significant'' windows are
recoverable from pure noise}

The same local-score procedure was run on the authoritative observed full-SNP scans already in
this pipeline (PC1, PC2, two bioclimatic covariates (bio6, bio11), bio\_winter, and the real
mitotype contrast), and each observed window was checked for physical overlap (same chromosome,
overlapping coordinate range) against the 20 same-mode null-replicate windows generated above (20
continuous or 20 mitoC2 null runs, as appropriate to the observed covariate being checked).

\begin{table}[H]
\centering\small
\begin{tabular}{lrrr}
\toprule
observed covariate & significant windows & overlap $\geq$1 null replicate & fraction \\
\midrule
PC1        & 7  & 3  & 42.9\% \\
PC2        & 6  & 2  & 33.3\% \\
bio6       & 7  & 2  & 28.6\% \\
bio11      & 6  & 0  & 0.0\% \\
bio\_winter & 7  & 2  & 28.6\% \\
\midrule
\emph{continuous, combined} & 33 & 9 & 27.3\% \\
\midrule
mitoC2     & 95 & 49 & 51.6\% \\
\bottomrule
\end{tabular}
\caption{Observed local-score windows that physically overlap at least one window independently
generated by a same-mode null replicate carrying no real signal.}
\end{table}

Roughly a quarter to a third of the observed continuous covariates' ``significant'' windows, and
just over half of the observed mitotype contrast's ``significant'' windows, occupy the same
physical genome locations as windows produced by pure-noise null replicates. Physical overlap does
not prove any individual observed window is itself a false positive -- a real signal and a null
artifact can coincidentally occupy the same low-recombination, LD-dense region, which is precisely
where both real and spurious local-score signal are known to concentrate. But it does mean that,
for a large fraction of the pipeline's reported ``significant'' local-score regions, \emph{the
observed result alone cannot distinguish a real association from a false positive that a
completely uninformative covariate would have produced anyway} at the same location.

\section*{\color{Accent}Answering the motivating question: is LD-aggregation necessary?}
No -- not to establish that \texttt{compute.local.scores()}'s raw significance output cannot be
trusted for this dataset. A modest, directly interpretable design -- 10 replicates, full SNP
resolution, no clustering or marker reduction, $\sim$2 hours of BayPass compute per run -- already
shows unambiguously that:
\begin{itemize}
\item a covariate with zero real signal produces well over 100 ``significant'' windows per
  replicate in the mitoC2/contrast mode, an order of magnitude more than in continuous mode, under
  an identical nominal threshold;
\item this failure holds for both structured and unstructured nulls, so it is not solely an
  artifact of population-structure correlation;
\item roughly a third to half of the pipeline's actual reported ``significant'' windows sit on top
  of locations where a null replicate independently produced a false positive.
\end{itemize}
Those three facts are sufficient, on their own, to conclude that local-score windows must be
interpreted through an explicit null comparison rather than at face value -- which is exactly what
this document set out to test, and did not require the LD-aggregated, large-$N$ (Stage-1/Stage-2,
$10^3$--$10^4$-replicate) machinery to show.

That said, this result does \emph{not} make the LD-aggregation approach pointless for other,
different purposes it was designed for: (i) building an actual calibrated empirical threshold or
FDR correction (this document establishes that raw significance is miscalibrated, but does not
itself produce a corrected threshold); (ii) explaining \emph{why} false positives cluster where
they do -- the companion Stage-2 LD-cluster redundancy analysis addresses that mechanistic
question directly, which a 10-replicate full-SNP design is not suited to answer at the
resolution needed. The claim here is narrower and specific to the question posed: LD-aggregation
is not \emph{necessary} to demonstrate the calibration failure itself.

\section*{\color{Accent}Limitations}
\begin{itemize}
\item $n=10$ replicates per null type is enough to see the calibration failure's magnitude clearly
  but is under-powered to resolve the structured-vs-unstructured contrast in mitoC2 mode at
  conventional significance ($p\approx0.06$--$0.08$); that comparison should be read as suggestive,
  not concluded.
\item Physical window overlap is a necessary but not sufficient condition for concluding an
  observed window is itself a false positive; it establishes that the location is null-susceptible,
  not that the specific observed association is spurious.
\item This analysis addresses only whether local-score's raw output is trustworthy at face value
  and whether a small design suffices to show that; it does not attempt to build a corrected
  significance threshold, is not conditioned on Stage-2 LD-cluster size or recombination rate, and
  does not draw on the Stage-1 1{,}000-replicate calibration (a separate, not-yet-run strand of
  this module).
\end{itemize}

\section*{\color{Accent}Reproducibility}
All 40 BayPass runs, the null covariate design, the local-score computation, and the
observed-vs-null overlap check are implemented fresh in this module
(\texttt{R/01\_generate\_null\_covariates.R}, \texttt{R/03\_compute\_local\_scores\_fullsnp.R},
\texttt{R/04\_observed\_vs\_null\_overlap.R}), reusing only the pipeline's authoritative inputs
(population covariance matrix, genotype matrix, observed BayPass scans) -- none of this document's
numbers are copied from, or dependent on, the earlier exploratory investigation of this question
(\texttt{module\_localscore\_crosscheck\_exploration/}, kept as read-only reference material).

\end{document}
